\documentclass[11pt]{article}
\usepackage[margin=1.1in]{geometry}
\usepackage{amsmath,amssymb,amsthm}
\usepackage[colorlinks=true,linkcolor=blue,citecolor=blue]{hyperref}

\newtheorem{theorem}{Theorem}[section]
\newtheorem{lemma}[theorem]{Lemma}
\newtheorem{corollary}[theorem]{Corollary}
\newtheorem{remark}[theorem]{Remark}

\newcommand{\R}{\mathbb{R}}
\newcommand{\E}{\mathbb{E}}
\newcommand{\PP}{\mathbb{P}}
\newcommand{\emp}{\bar\mu_N}
\newcommand{\vol}{\operatorname{vol}}

\title{The interpolation law for the critical scaling of the\\
uniform doubling condition}
\author{Yitzchak Shmalo\\
\small Einstein Institute of Mathematics, The Hebrew University of Jerusalem, Jerusalem, Israel\\
\small \texttt{yitzchak.shmalo@gmail.com}}
\date{June 2026}

\begin{document}
\maketitle

\begin{abstract}
This paper completes the program, carried out in six companion papers, of
determining the critical scaling of the uniform doubling condition of
[Shmalo, \emph{Ann.\ Math.\ Artif.\ Intell.} (2024)] for independent samples.
For the model truncation regions
$G_k=B^n\cap\{x_1\le0,\dots,x_k\le0\}$, $0\le k\le n-2$, whose lowest
transversal boundary stratum --- the sphere
$\{x\in S^{n-1}:x_1=\dots=x_k=0\}$ --- has dimension $j=n-1-k$, we prove that
the critical scaling factor is exactly
\[
\ell_N\;\asymp\;\Bigl(\frac{\log N}{N}\Bigr)^{2/(2n-j)}
\;=\;\Bigl(\frac{\log N}{N}\Bigr)^{2/(n+1+k)},
\]
with both the upper bound (the condition holds with probability $1-1/N$ at
this scaling) and the lower bound (it fails with probability $1-o(1)$ below
this scaling, logarithm included) established for every $k\le n-2$; for
$k\ge n-1$ the conical exponent $\frac1n$ of the companion papers applies.
This proves the interpolation law conjectured in
[\emph{The edge exponent of the uniform doubling condition}] at every stratum
dimension, in the orthogonal model.  The upper bound rests on a mixed-regime
section lemma: at depth $h$ below the extreme point of $G_k$ in an arbitrary
direction, the level section has volume $\ge c\,h^{(n-1+k)/2}$, proved by
splitting the truncated coordinates at the angular threshold
$\sqrt h$ --- steep coordinates contribute a factor $h$ through a fibre box,
shallow ones are absorbed into a spherical cap centered on the stratum, of
which the all-negative orthant keeps the exact fraction $2^{-k}$ by reflection
symmetry.  The lower bound is a second-moment argument over
$\asymp(N/\log N)^{j/(2n-j)}$ congruent caps along the stratum.
\end{abstract}

\section{Setting and main theorem}\label{sec:intro}

Let $n\ge2$, let $\mu$ be the uniform probability measure on the closed unit
ball $K=B^n$, and let $T=\{X_1,\dots,X_N\}$ be i.i.d.\ samples with empirical
measure $\emp$.  For $0\le k\le n-2$ consider the truncation set
$h_k:=\{(e_1,0),\dots,(e_k,0)\}$ with region
\[
G\;=\;G_k\;:=\;\bigl\{x\in B^n:\ x_1\le0,\dots,x_k\le0\bigr\},
\]
and the uniform doubling condition (UDC) of \cite{ShmaloUDC} for slabs
truncated by $h_k$, with constants $(\kappa,\delta,\sigma,\beta,\ell)$, as in
the companion papers \cite{companion2,companion3,companion4,companion5,
companion6,companion7}.  The boundary of $G_k$ is stratified by the subsets
$S\subseteq\{1,\dots,k\}$ of active constraints; its lowest-dimensional
spherical stratum is
\[
E\;:=\;\bigl\{x\in S^{n-1}:\ x_1=\dots=x_k=0\bigr\}\;\cong\;S^{\,j},
\qquad
j\;:=\;n-1-k\;\ge\;1 .
\]
The region is uniformly thick: with
$c^-:=-\frac1{2\sqrt k}(e_1+\dots+e_k)$ (for $k\ge1$; $c^-:=0$ for $k=0$),
\begin{equation}\label{eq:thick}
B\Bigl(c^-,\ \tfrac1{2\sqrt k}\wedge\tfrac12\Bigr)\;\subseteq\;G_k .
\end{equation}

\begin{theorem}[Interpolation law, orthogonal model]\label{thm:main}
Fix $n\ge2$ and $0\le k\le n-2$, and fix $\kappa>1$, $\sigma>0$, $\beta>0$,
$\delta\ge\frac3N$ where relevant.
\begin{enumerate}
\item[(i)] {\rm(Upper bound.)}  With the $N$-independent constants of
\cite{companion2} {\rm(}$\kappa=151$, $\sigma=1$, $\delta=\delta(n,k)>0$,
$\beta=2$, thickness radius $r_0=\frac1{2\sqrt k}\wedge\frac12${\rm)}, there
are $C(n,k)$ and $N_0$ such that for $N\ge N_0$, with probability at least
$1-\frac1N$, the sample satisfies the UDC for the truncation set $h_k$ with
scaling factor
\[
\ell_N\;=\;C(n,k)\,\Bigl(\frac{\log N}{N}\Bigr)^{2/(n+1+k)} .
\]
\item[(ii)] {\rm(Lower bound.)}  There are $c(n,k,\kappa)>0$ such that with
probability $1-o(1)$, the doubling inequality fails for some slab truncated by
$h_k$ at every grid scale below $c\,(\log N/N)^{2/(n+1+k)}$; hence the UDC
requires $\ell\ge c\,(\log N/N)^{2/(n+1+k)}$.
\end{enumerate}
Consequently the critical scaling is exactly
$\Theta\bigl((\log N/N)^{2/(2n-j)}\bigr)$ for every stratum dimension
$1\le j\le n-1$, while for truncation geometry with conical points
{\rm(}$j=0${\rm)} it is $\Theta\bigl((\log N/N)^{1/n}\bigr)$ by
\cite{companion3,companion5}.  The interpolation law of
\cite[Conjecture 1.2]{companion6} holds at every $j$ in the orthogonal model.
\end{theorem}

The cases $k=0,1$ are the companion results
\cite{companion4,companion6,companion7}; we treat all $k\le n-2$ uniformly,
and the reader may keep $k=2$, $n=4$ in mind.  Throughout,
$x=(x_1,\dots,x_k,x^\flat)$ splits coordinates into the truncated ones and the
rest, and for a unit vector $u$ we write $u=(u_1,\dots,u_k,u^\flat)$.

\section{Caps at the stratum}\label{sec:caps}

\begin{lemma}[Extreme points]\label{lem:extreme}
Let $u$ be a unit vector, $P:=\{i\le k:\ u_i>0\}$, and let
$\hat u$ be $u$ with the $P$-coordinates replaced by $0$.  If
$\hat u\ne0$, then
\[
\max_{x\in G_k}\langle x,u\rangle\;=\;|\hat u|,
\]
attained exactly at $p:=\hat u/|\hat u|\in G_k$, which satisfies $p_i=0$ for
$i\in P$ and $p_i\le0$ for $i\le k$, $i\notin P$.
\end{lemma}

\begin{proof}
For $x\in G_k$: $\langle x,u\rangle=\sum_{i\in P}u_ix_i
+\langle x,\hat u\rangle\le0+|x|\,|\hat u|\le|\hat u|$, since $u_i>0$ and
$x_i\le0$ for $i\in P$; equality forces $x_i=0$ ($i\in P$), $|x|=1$ and
$x\parallel\hat u$.  Finally $p\in G_k$ because $\hat u_i=u_i\le0$ for
$i\le k$, $i\notin P$.
\end{proof}

If $\hat u=0$ (i.e.\ $u$ is supported on $P$), the maximum $0$ is attained on
the entire face stratum $\{x\in G_k:\ x_i=0,\ i\in P\}$, a set of positive
$(n-|P|)$-dimensional volume; this degenerate case only makes the section
bounds below easier and is handled separately in Lemma~\ref{lem:section}.

For $p\in E$ define the \emph{stratum direction} and \emph{stratum cap}
\[
u_p:=\frac{e_1+\dots+e_k+p}{\sqrt{k+1}},
\qquad
C_p(h):=\Bigl\{x\in G_k:\ \langle x,u_p\rangle\ge\tfrac1{\sqrt{k+1}}-h\Bigr\}.
\]
By Lemma~\ref{lem:extreme} (here $P=\{1,\dots,k\}$, $\hat u_p=p/\sqrt{k+1}$),
the maximum of $\langle x,u_p\rangle$ over $G_k$ is $\frac1{\sqrt{k+1}}$,
attained at $p$.

\begin{lemma}[Stratum cap volume and footprint]\label{lem:capmass}
There are $0<c_1(n,k)\le c_2(n,k)$ and $h_1(n,k)>0$ such that for all
$p\in E$ and $0<h\le h_1$,
\[
c_1\,h^{\frac{n+1+k}2}\;\le\;\mu\bigl(C_p(h)\bigr)\;\le\;
c_2\,h^{\frac{n+1+k}2},
\qquad
C_p(h)\subseteq\{x:\ |x-p|\le c_3(k)\sqrt h\,\}.
\]
\end{lemma}

\begin{proof}
By the symmetry of $G_k$ under rotations fixing $e_1,\dots,e_k$ we may take
$p=e_{k+1}$.  Write $x_i=-t_i$ ($t_i\ge0$) for $i\le k$, let
$y:=x_{k+1}$, and let $x''\in\R^{n-k-1}$ collect the remaining coordinates.
With $h':=h\sqrt{k+1}$ the cap is
\[
\Bigl\{\,t_i\ge0,\ |x|\le1,\ y\;\ge\;1-h'+\textstyle\sum_it_i\,\Bigr\},
\]
and $|x|\le1$ gives $y\le\sqrt{1-\sum t_i^2-|x''|^2}\le1-\frac{\sum
t_i^2+|x''|^2}2$.  \emph{Upper bound:} the constraints force
$\sum_it_i\le h'$ (so each $t_i\le h'$), $|x''|^2\le2h'$, and the
$y$-fibre has length at most $h'$; hence
$\vol\le h'\cdot(h')^k\cdot\omega_{n-k-1}(2h')^{\frac{n-k-1}2}
\le C h^{\frac{n+1+k}2}$.  \emph{Lower bound:} restrict to
$t_i\le\frac{h'}{8k}$ ($i\le k$) and $|x''|^2\le\frac{h'}2$; then
$\sum t_i\le\frac{h'}8$, $\sum t_i^2\le\frac{h'^2}{64k}\le\frac{h'}{64}$
(for $h'\le1$), and the fibre length is at least
\[
\Bigl(1-\tfrac{\sum t_i^2+|x''|^2}2\Bigr)-\Bigl(1-h'+\sum t_i\Bigr)
\;\ge\;h'\Bigl(1-\tfrac18-\tfrac1{128}-\tfrac14\Bigr)\;\ge\;\tfrac{h'}2,
\]
using $\sum t_i\le\frac{h'}8$, $\frac{\sum t_i^2}2\le\frac{h'}{128}$ and
$\frac{|x''|^2}2\le\frac{h'}4$; so
$\vol\ge\frac{h'}2(\frac{h'}{8k})^k\omega_{n-k-1}(\frac{h'}2)^{\frac{n-k-1}2}
\ge c\,h^{\frac{n+1+k}2}$.  \emph{Footprint:} on the cap,
$1\ge y\ge1-h'$, $t_i\le h'$, $|x''|^2\le2h'$, so
$|x-p|^2\le kh'^2+h'^2+2h'\le c_3^2h$.
\end{proof}

\begin{lemma}[Stall]\label{lem:stall}
Fix $\kappa>1$, $\sigma>0$, $\delta\ge\frac3N$, $p\in E$ and
$0<h\le\frac H4$, $H\le h_1$.  If exactly one sample point lies in
$C_p(H)$ and it lies in $C_p(h)$, then for every $w\le H/\kappa$ the slab
with direction $u_p$, center that sample point, and truncation set $h_k$
violates the doubling inequality; hence the UDC fails for every
$\ell\le H/\kappa^2$.
\end{lemma}

\begin{proof}
Identical to the companion stall lemmas: for $w'\le H$ the slab is contained
in $C_p(h+\frac{w'}2)\subseteq C_p(H)$, so the empirical masses of $S_w$ and
$S_{\kappa w}$ both equal $\frac1N<\min\{\delta,(1+\sigma)\frac1N\}$.
\end{proof}

\section{The lower bound}\label{sec:lower}

\begin{proof}[Proof of Theorem \ref{thm:main}(ii)]
This is the second-moment scheme of
\cite{companion4,companion5,companion6} verbatim, with the cap data of
Lemma~\ref{lem:capmass}.  Set $q:=(4\kappa^2)^{\frac{n+1+k}2}\frac{c_2}{c_1}$,
$\varepsilon_3:=\frac{j}{4q(2n-j)}$, $\tau_N:=\varepsilon_3\log N$; define
$h_N$ by $\mu(C_p(h_N))=\frac{\tau_N}N$ (independent of $p$ by symmetry), so
$h_N\asymp(\tau_N/N)^{2/(n+1+k)}$, and $H_N:=4\kappa^2h_N$, so
$\mu(C_p(H_N))\le\frac{q\tau_N}N$.  By Lemma~\ref{lem:capmass} the caps have
footprint $c_3\sqrt{H_N}$, so a maximal $2c_3\sqrt{H_N}$-separated family
$p_1,\dots,p_M\in E\cong S^{\,j}$ gives pairwise disjoint outer caps with
$M\ge c\,H_N^{-j/2}\ge c'(N/\tau_N)^{\frac{j}{n+1+k}}$ --- note
$\frac{j}{n+1+k}=\frac{j}{2n-j}$.  With
$E_\alpha$ the event that $C_{p_\alpha}(H_N)$ contains exactly one sample
point and that point lies in $C_{p_\alpha}(h_N)$, and
$X:=\sum_\alpha\mathbf1_{E_\alpha}$:
\[
p:=\PP(E_\alpha)=\tau_N\bigl(1-\mu(C(H_N))\bigr)^{N-1}
\ge\tau_NN^{-2q\varepsilon_3},
\qquad
\E X=Mp\ge c'N^{\frac j{2n-j}-2q\varepsilon_3}\tau_N^{1-\frac j{2n-j}}
\ge c'N^{\frac j{2(2n-j)}}\to\infty,
\]
and the disjoint-cap correlation ratio is
$\exp(O(N\mu(C(H_N))^2))=1+o(1)$, so
$\PP(X=0)\le(\E X)^{-1}+o(1)=o(1)$ by Chebyshev.  On $\{X\ge1\}$,
Lemma~\ref{lem:stall} forces
$\ell>H_N/\kappa^2=4h_N\ge c(\log N/N)^{2/(n+1+k)}$.
\end{proof}

\section{The section lemma}\label{sec:section}

For a unit vector $u$ with $\hat u\ne0$ let
$\tau(x):=\langle x,u\rangle$, $\tau_{\max}:=|\hat u|$ (Lemma
\ref{lem:extreme}), and
$A_u(h):=\vol_{n-1}(G_k\cap\{\tau=\tau_{\max}-h\})$.

\begin{lemma}[Section lemma]\label{lem:section}
There are $a(n,k)>0$ and $h_0(n,k)>0$ such that for every unit vector $u$ and
every $0<h\le h_0$,
\[
A_u(h)\;\ge\;a(n,k)\,h^{\frac{n-1+k}2} .
\]
\end{lemma}

\begin{proof}
As in \cite{companion7}, it suffices to prove the cap-mass bound
\begin{equation}\label{eq:cap}
\Lambda^{(n,k)}_u(h):=\mu\bigl(\{x\in G_k:\ \tau(x)\ge\tau_{\max}-h\}\bigr)
\;\ge\;b(n,k)\,h^{\frac{n+1+k}2}
\qquad(0<h\le h_0),
\end{equation}
for every unit $u$ with $\hat u\ne0$, together with a bulk anchor; the
conversion to sections is given at the end.  We prove \eqref{eq:cap} by
induction on $k$; the base case $k=0$ is the standard spherical cap bound
$\Lambda^{(n,0)}_u(h)\ge c(n)h^{\frac{n+1}2}$.

\emph{Case 1: $|\hat u|\ge\frac12$.}  With $p=\hat u/|\hat u|$, split the
truncated indices into
\[
\mathcal A:=\Bigl\{i\le k:\ u_i\ge\tfrac{\sqrt h}{2k}\Bigr\},
\quad
\mathcal Z:=\Bigl\{i\le k:\ i\notin\mathcal A,\
u_i>0\ \text{ or }\ |p_i|<2\sqrt h\Bigr\},
\quad
\mathcal N:=\{1,\dots,k\}\setminus(\mathcal A\cup\mathcal Z),
\]
so that $u_i\le0$ and $p_i\le-2\sqrt h$ for $i\in\mathcal N$.  Let
$\tilde u$ be $u$ with the $\mathcal A\cup\mathcal Z$-coordinates replaced by
$0$.  Every coordinate zeroed in $\tilde u$ but present in $\hat u$ lies in
$\mathcal Z\setminus P$ and satisfies
$|\hat u_i|=|u_i|\le\max\bigl(\tfrac{\sqrt h}{2k},\,2\sqrt h\bigr)\le2\sqrt h$,
whence
\[
|\hat u|^2-|\tilde u|^2\;\le\;4kh,
\qquad
|\hat u|-|\tilde u|\;\le\;\frac{4kh}{|\hat u|}\;\le\;8kh,
\]
and $|\tilde u|\in[\tfrac14,1]$ for $h\le h_0(k)$.  Set
$\tilde p:=\tilde u/|\tilde u|\in G_k$; then $\tilde p_i=0$ for
$i\in\mathcal A\cup\mathcal Z$, and for $i\in\mathcal N$,
$|\tilde p_i|=|u_i|/|\tilde u|\ge|u_i|\ge2\sqrt h\,|\hat u|\ge\sqrt h$, with
$\tilde p_i<0$.

Decompose $x=y\,\tilde p+w$, $w\perp\tilde p$.  Since
$\tilde p\perp e_i$ for $i\in\mathcal A\cup\mathcal Z$, the coordinates
$w_i:=x_i$ ($i\in\mathcal A\cup\mathcal Z$) are coordinates of the
hyperplane $\tilde p^\perp$; let $w^\perp$ denote the remaining component of
$w$.  Define the witness box $W$ by
\[
w_i\in\Bigl[-\min\Bigl(\frac{h}{8k\,u_i},\frac{\sqrt h}8\Bigr),\,0\Bigr]
\ (i\in\mathcal A),
\qquad
w_i\in\Bigl[-\frac{\sqrt h}{8\sqrt k},\,0\Bigr]\ (i\in\mathcal Z),
\qquad
|w^\perp|\le\frac{\sqrt h}8,
\]
\[
1-\frac{h}{4|\tilde u|}\;\le\;y\;\le\;\sqrt{1-|w|^2}.
\]

\emph{$W$ is nonempty with a long fibre.}  We have
$|w|^2\le k\cdot\frac h{64k}+k\cdot\frac h{64k}+\frac h{64}\le\frac{3h}{64}$,
so $\sqrt{1-|w|^2}\ge1-\frac{3h}{64}$, while
$\frac h{4|\tilde u|}\ge\frac h4$; the $y$-fibre has length at least
$\frac h4-\frac{3h}{64}\ge\frac h8$.

\emph{$W\subseteq G_k$.}  For $i\in\mathcal A\cup\mathcal Z$:
$x_i=w_i\le0$.  For $i\in\mathcal N$:
$x_i=y\tilde p_i+(w^\perp)_i\le(1-h)(-\sqrt h)+\frac{\sqrt h}8<0$ for
$h\le\frac12$.  And $|x|^2=y^2+|w|^2\le1$.

\emph{$W$ lies in the cap, after rescaling.}  Since $w\perp\tilde p$ and
$u-\tilde u$ is supported on $\mathcal A\cup\mathcal Z$,
\[
\tau(x)=y\,\langle\tilde p,u\rangle+\langle w,u\rangle
=y\,|\tilde u|+\sum_{i\in\mathcal A\cup\mathcal Z}u_iw_i ,
\]
using $\langle\tilde p,u\rangle=\langle\tilde p,\tilde u\rangle=|\tilde u|$
and $\langle w,\tilde u\rangle=|\tilde u|\langle w,\tilde p\rangle=0$.  Now
$\sum_{\mathcal A}u_i|w_i|\le k\cdot\frac h{8k}=\frac h8$, and on
$\mathcal Z$ every coordinate has $|u_i|\le2\sqrt h$ {\rm(}either
$0<u_i<\frac{\sqrt h}{2k}$, or $u_i\le0$ with
$|u_i|=|p_i|\,|\hat u|<2\sqrt h${\rm)}, so
$\sum_{\mathcal Z}|u_i||w_i|\le k\cdot2\sqrt h\cdot
\frac{\sqrt h}{8\sqrt k}\le\frac{\sqrt k\,h}{4}$.  Hence on $W$,
\[
\tau(x)\;\ge\;\Bigl(1-\frac h{4|\tilde u|}\Bigr)|\tilde u|
-\frac h8-\frac{\sqrt k\,h}{4}
\;\ge\;|\tilde u|-\Bigl(\frac38+\frac{\sqrt k}4\Bigr)h
\;\ge\;|\hat u|-8kh-kh
\;\ge\;\tau_{\max}-9kh ,
\]
using $\frac38+\frac{\sqrt k}4\le k$ for $k\ge1$.
Replacing $h$ by $h/(9k)$ in the whole construction (which only changes
constants), $W$ lies in the cap of depth $h$.

\emph{Volume.}  Each $\mathcal A$-window has length
$\min\bigl(\frac h{8ku_i},\frac{\sqrt h}8\bigr)\ge\frac h{8k}$ (using
$u_i\le1$ and $h\le k^2$); each $\mathcal Z$-window has length
$\frac{\sqrt h}{8\sqrt k}$; the $w^\perp$-ball has dimension
$n-1-|\mathcal A|-|\mathcal Z|$; the fibre has length $\ge\frac h8$.  Hence
\[
\Lambda^{(n,k)}_u(h)\;\ge\;\frac{\vol(W)}{\vol(B^n)}
\;\ge\;c(n,k)\,h\cdot h^{|\mathcal A|}\cdot
h^{\frac{|\mathcal Z|+(n-1-|\mathcal A|-|\mathcal Z|)}2}
\;=\;c(n,k)\,h^{\frac{n+1+|\mathcal A|}2}
\;\ge\;c(n,k)\,h^{\frac{n+1+k}2},
\]
proving \eqref{eq:cap} in Case 1.

\emph{Case 2: $|\hat u|<\frac12$ (induction on $k$).}  Then
$\sum_{i\in P}u_i^2=1-|\hat u|^2>\frac34$, so some $i_0\in P$ has
$u_{i_0}>\frac{\sqrt3}{2\sqrt k}\ge\frac1{2\sqrt k}$.  Write
$x=(x_{i_0},x')$ and $u=(u_{i_0},u')$; note $u'\ne0$ would fail only if
$u=e_{i_0}$, in which case $\hat u=0$, excluded.  For
$s\in[-\frac{h}{2u_{i_0}},0]$ --- an interval of length
$\frac h{2u_{i_0}}\ge\frac h2$ --- the slice
$G_k\cap\{x_{i_0}=s\}$ contains $\sqrt{1-s^2}\,\cdot G'$, where
$G'$ is the orthogonal model with parameters $(n-1,k-1)$ in the coordinates
$x'$ {\rm(}the slice is exactly the ball of radius $\sqrt{1-s^2}$
intersected with the $k-1$ remaining half-spaces{\rm)}.  On the slice,
$\tau(x)=u_{i_0}s+\langle x',u'\rangle\ge-\frac h2+\langle x',u'\rangle$.
The maximum of $\langle x',u'\rangle$ over the unscaled slice region equals
$|\widehat{u'}|=|\hat u|$ {\rm(}zeroing the positive truncated coordinates of
$u'$ gives the same vector $\hat u$, supported away from
$i_0${\rm)}, and scaling by $\sqrt{1-s^2}\ge1-h$ changes it by at most
$|\hat u|h\le h$.  Hence the cap of depth $h$ for $(G_k,u)$ contains, over
each slice, the image of the cap of depth $h_*$ for
$(G',u'/|u'|)$ with
$h_*:=\frac{h/4}{|u'|}\ge\frac h4$, provided
$u_{i_0}|s|+\,$scaling loss$\,\le\frac h2+h\le2h$ and the slice-cap depth
$\frac h4$ are budgeted within a total depth $4h$; replacing $h$ by
$h/4$ throughout, the inclusion holds at depth $h$.  By the induction
hypothesis applied to $(n-1,k-1)$ at depth $\min(h_*,h_0)$ --- if
$h_*>h_0$ use monotonicity $\Lambda(h_*)\ge\Lambda(h_0)\ge b\,h_0^{(n+k-1)/2}
\ge b'h^{(n+k-1)/2}$ --- and Fubini over $s$,
\[
\Lambda^{(n,k)}_u(h)\;\ge\;c\cdot\frac h2\cdot
b(n-1,k-1)\,h^{\frac{(n-1)+1+(k-1)}2}
\;=\;c'\,h^{1+\frac{n+k-1}2}\;=\;c'\,h^{\frac{n+1+k}2}.
\]

\emph{Conversion to sections.}  Exactly as in \cite{companion7}: by
Brunn--Minkowski, $A_u$ is unimodal in the depth; the inner ball
\eqref{eq:thick} provides a bulk depth
$h^\dagger\ge c(k)>0$ with $A_u(h^\dagger)\ge a_1(n,k)>0$ {\rm(}its
$\tau$-range is an interval of length $2r_0$ inside the range of $G_k$,
whose sections at distance $\le\frac{r_0}2$ from its center have radius
$\ge\frac{\sqrt3}2r_0${\rm)}.  For $h\le h_0\le h^\dagger$: if $h$ is below
the mode of $A_u$, then $A_u$ is nondecreasing on $[0,h]$ and
$A_u(h)\ge\Lambda_u(h)/h\ge b\,h^{\frac{n-1+k}2}$ by \eqref{eq:cap}; if $h$
is above the mode, unimodality gives
$A_u(h)\ge A_u(h^\dagger)\ge a_1\ge a_1h^{\frac{n-1+k}2}$.
\end{proof}

\begin{remark}
The witness box exploits the key orthogonality $\langle w,\tilde u\rangle=0$:
once the shallow coordinates are zeroed into the anchor direction
$\tilde p$, the level constraint involves only the $\mathcal A\cup\mathcal
Z$ coordinates, so all remaining directions are free at scale $\sqrt h$ and
the volume count telescopes.  The induction of Case 2 is the mechanism by
which face strata inherit the law from lower-dimensional models, one
truncation per dimension.
\end{remark}

\section{Slab bound and proof of the upper bound}\label{sec:upper}

\begin{theorem}[Slab volume]\label{thm:slab}
There are $c(n,k)>0$, $L_0(n,k)>0$ with
\[
\vol\Bigl(G_k\cap\{|\langle y-x,u\rangle|\le L\}\Bigr)
\;\ge\;c\,L^{\frac{n+1+k}2}
\qquad(x\in G_k,\ u\in S^{n-1},\ 0<L\le L_0).
\]
\end{theorem}

\begin{proof}
Anchor--unimodality as in \cite{companion3,companion7}: anchors at depth
$\frac L4$ from each $u$-extreme with value
$a(n,k)(\frac L4)^{\frac{n-1+k}2}$ by Lemma~\ref{lem:section} applied to
$u$ and $-u$; unimodality spreads the minimum anchor value over the middle
of the range, whose complement has length $\frac L2$; the range has length at
least $2r_0\ge2L$ by \eqref{eq:thick}; and the slab retains length
$\ge L$ in the range, of which $\ge\frac L2$ in the middle.
\end{proof}

\begin{proof}[Proof of Theorem \ref{thm:main}(i)]
Insert Theorem~\ref{thm:slab} into the proof skeleton of
\cite{companion2,companion3}: the population dichotomy holds with
$r_0=\frac1{2\sqrt k}\wedge\frac12$ by \eqref{eq:thick}; the VC relative
deviation bounds apply to slabs with $k$ truncations; and a sample point in
$S_w$ yields $\mu(S_{\kappa w})\ge32u_N$ for all
$w\ge\ell_N:=C(n,k)(u_N)^{2/(n+1+k)}\asymp(\log N/N)^{2/(n+1+k)}$.
\end{proof}

\section{Concluding remarks}\label{sec:concl}

The orthogonal model is the normal form of transversal stratified truncation:
for half-spaces in general position meeting $\partial K$ and each other at
dihedral angles bounded below, the same two-step scheme applies with
constants depending on the angles, the witness-box and reflection arguments
being replaced by their affine images.  The UDC constants themselves are stable
under a fixed nonsingular affine change of variables: by the affine-transfer
lemma of the high-probability companion paper \cite{companion2}, an affine map
with singular values in $[\lambda,\Lambda]$ sends constants
$(\kappa,\delta,\sigma,\beta,\ell)$ to
$(\kappa^2,\delta,\sigma,\lambda\beta/\kappa,\Lambda\ell)$ in the grid form,
and to $(\kappa,\delta,\sigma,\lambda\beta,\Lambda\ell)$ in the all-width form.
Thus extending the interpolation law beyond the orthogonal model reduces to the
local section-volume estimates, not to a new doubling-formalism issue.  Within
the model, the program of the seven papers is complete:
for every stratum dimension $j\in\{0,1,\dots,n-1\}$ the critical scaling of
the uniform doubling condition is
\[
\ell_N\;\asymp\;\Bigl(\frac{\log N}{N}\Bigr)^{2/(2n-j)}
\quad(1\le j\le n-1),
\qquad
\ell_N\;\asymp\;\Bigl(\frac{\log N}{N}\Bigr)^{1/n}
\quad(j=0),
\]
exponents and logarithms exact.  In the language of the stability theory of
\cite{BJS,ShmaloUDC}: the finest resolution at which i.i.d.\ data can certify
the doubling condition is set by the most isolated extremal sample point of
the worst boundary stratum of the activation geometry, and degrades
gracefully --- through the explicit exponent ladder
$\frac{2}{n+1},\frac{2}{n+2},\dots,\frac1n$ --- as the truncation geometry
sharpens from smooth to conical.

\subsection*{Acknowledgments}
The research presented in this paper was supported by the European Research
Council (ERC) under the European Union's Horizon Europe research and
innovation programme (grant agreement No.\ 101041711), by the Israel Science
Foundation (grant number 2258/19) and by the Israel Science Foundation (ISF
Grant 4101/25).

\subsection*{Declaration on generative AI}
During the preparation of this work the author used generative AI tools to
accelerate drafting and revision.  All mathematical claims, proofs, and
bibliographic information were subsequently reviewed and edited by the
author, who takes full responsibility for the content of the manuscript.

\begin{thebibliography}{9}

\bibitem{BJS}
L.~Berlyand, P.-E.~Jabin and C.~A.~Safsten,
\emph{Stability for the training of deep neural networks and other
classifiers}, Math.\ Models Methods Appl.\ Sci.\ \textbf{31} (2021),
2345--2390.

\bibitem{ShmaloUDC}
Y.~Shmalo,
\emph{Stability of accuracy for the training of DNNs via the uniform doubling
condition}, Ann.\ Math.\ Artif.\ Intell.\ (2024),
DOI 10.1007/s10472-023-09919-1.

\bibitem{companion2}
Y.~Shmalo, \emph{The uniform doubling condition holds with high probability
for independent samples}, manuscript, 2026.

\bibitem{companion3}
Y.~Shmalo, \emph{The critical scaling of the uniform doubling condition:
smooth bodies versus truncations}, manuscript, 2026.

\bibitem{companion4}
Y.~Shmalo, \emph{The logarithm is necessary: exact critical scaling of the
uniform doubling condition on the ball}, manuscript, 2026.

\bibitem{companion5}
Y.~Shmalo, \emph{The exact critical scaling of the uniform doubling
condition with truncations}, manuscript, 2026.

\bibitem{companion6}
Y.~Shmalo, \emph{The edge exponent of the uniform doubling condition and an
interpolation conjecture}, manuscript, 2026.

\bibitem{companion7}
Y.~Shmalo, \emph{The edge exponent of the uniform doubling condition, II:
the matching upper bound}, manuscript, 2026.

\end{thebibliography}

\end{document}
